\documentclass[12pt]{extarticle}
\def\activityid{F04-FAC-v2}\usepackage[letterpaper,margin=.65in,top=.60in,bottom=.65in]{geometry}
\usepackage[T1]{fontenc}\usepackage[default]{sourcesanspro}
\usepackage{microtype,tikz,array,tabularx,fancyhdr,amsmath,amssymb}
\usepackage[hidelinks]{hyperref}\usetikzlibrary{calc,arrows.meta}
\setlength{\parindent}{0pt}\setlength{\parskip}{7pt}\setlength{\headheight}{0pt}\setlength{\footskip}{26pt}
\setlength{\emergencystretch}{2em}\pagestyle{fancy}\fancyhf{}\renewcommand{\headrulewidth}{0pt}\renewcommand{\footrulewidth}{.4pt}
\fancyfoot[L]{\scriptsize Bellingham Math Circle / Week 4 / \activityid}\fancyfoot[R]{\scriptsize \thepage}
\definecolor{ink}{gray}{.12}\definecolor{soft}{gray}{.95}\color{ink}
\newcommand{\PageTitle}[2]{{\fontsize{10}{12}\selectfont\bfseries WEEK 4\quad ROUND AND ROUND\hfill #1\par}\vspace{3pt}{\fontsize{26}{29}\selectfont\bfseries #2\par}\vspace{2pt}}
\newcommand{\NameLine}{{\small Name: \rule{2.65in}{.4pt}\hfill Date: \rule{1.35in}{.4pt}\par}\vspace{4pt}}
\newcommand{\Task}[2]{\par\vspace{3pt}{\large\bfseries #1}\par #2\par}
\newcommand{\Subhead}[1]{\par\vspace{3pt}{\large\bfseries #1}\par}
\newcommand{\RuleBox}[1]{\par\begingroup\setlength{\fboxsep}{8pt}\colorbox{soft}{\parbox{\dimexpr\linewidth-16pt\relax}{#1}}\endgroup\par}
\newcommand{\WriteBox}[2]{\par\textbf{#1}\par\vspace{2pt}\begin{tikzpicture}\draw[black!40,line width=.5pt] (0,0) rectangle (\dimexpr\linewidth-.5pt\relax,#2);\end{tikzpicture}\par}
\newcommand{\StudentFooter}{\vfill{\small Try it with objects. Explain by pointing, talking, or drawing.}\par}
\newcommand{\LampRing}[3]{\begin{tikzpicture}[x=1in,y=1in]
\foreach \i in {1,...,#1}{\coordinate (v\i) at ({90-360*(\i-1)/#1}:#2);}
\foreach \i in {1,...,#1}{\pgfmathtruncatemacro{\j}{mod(\i,#1)+1}\draw[thick] (v\i)--node[midway,fill=white,inner sep=2pt,font=\small]{\i\j}(v\j);}
\foreach \i in {1,...,#1}{\node[circle,draw,fill=white,minimum size=.52in,inner sep=0pt,font=\bfseries] at(v\i){\i};}
\foreach \i in {#3}{\node[circle,draw,fill=ink,text=white,minimum size=.52in,inner sep=0pt,font=\bfseries] at(v\i){\i};}
\end{tikzpicture}}
\newcommand{\Lights}[1]{\begin{tikzpicture}[x=.55in,y=.55in,baseline=-.10in]\foreach \v [count=\i] in {#1}{\ifnum\v=1\fill (\i,0) circle(.16in);\else\draw (\i,0) circle(.16in);\fi}\end{tikzpicture}}
\newcommand{\ClockRing}[2]{\begin{tikzpicture}[x=1in,y=1in]\draw[black!30] (0,0) circle(#2);\pgfmathtruncatemacro{\n}{#1-1}\foreach\i in {0,...,\n}{\node[circle,draw,fill=white,minimum size=.35in,inner sep=0pt] at({90-360*\i/#1}:#2){\i};}\draw[-{Stealth},thick] (-.35,.15) arc(155:25:.38);\node[font=\small] at(0,-.18){clockwise};\end{tikzpicture}}
\newcommand{\Key}[2]{\begin{center}\renewcommand{\arraystretch}{1.55}\begin{tabular}{c|*{#1}{c}}in&#2\end{tabular}\end{center}}

\begin{document}

\PageTitle{FACILITATOR / 1}{Rotations, orbits, and shuffles}
\RuleBox{Week 3 allowed arbitrary permutations. Week 4 fixes a circular order and studies shifts: what controls their cycles and return times? The extra changes addition to multiplication through a real card shuffle.}
\Subhead{Materials and preparation}
Print 3 K--1, 3 middle, 1 upper, 1 extra reserve, and this guide. Bring seven counters, pencils, scratch paper, and cards 0--7 (another eight for the sixteen-card extension). For the K six-picture extension, draw shapes on scraps. Hop counters on the printed fixed rings; no movable wheel or brad is needed. Draw strips on scratch paper for the optional proof.
\Subhead{Readiness}
\textbf{K--1:} match shapes, count hops to three with help; no independent reading. Reason about repeated movement and impossible landings. \textbf{2--3:} A--J familiarity, count up to ten; read-aloud removes word-reading barriers. Explain disjoint routes and use a common key. \textbf{4--5:} count to eighteen for finite rings. The optional general proof additionally needs division/gcd and the adult-supported coprime lemma below. Distinguish an orbit from the whole permutation. \textbf{Extra:} track numbered positions and remainders; powers can be explained as repeated doubling. The core eight-card proof needs only 2,4,8 and remainder 1.
\Subhead{A shared 60-minute hour}
\textbf{0--10:} play with counters and rings, trying hop lengths. \textbf{10--15:} define clockwise and count movements, not the starting place. Launch: ``Can a step rule miss some places forever?'' \textbf{15--35:} K compares one/two-step routes; middle tries the code clue; upper explores the 12-ring. Adult A anchors K; organizer visits middle and upper. \textbf{35--40:} movement/reset: stand, stretch, leave work ready to return to. \textbf{40--55:} continue orbit comparisons or an optional proof below; offer the shuffle as an alternative. \textbf{55--60:} share a route, obstruction, or explanation; tidy counters and cards.
The twelve-row table is optional: inverse shifts share return times. Move to explanation when ready. An adult should discuss the upper child's conjecture; the three pages are choices, not a timed test.
\Subhead{Hints and common confusions}
Count spaces moved, not the starting label. Stop an orbit when its start returns; then choose an unused start. Ask how many spaces and whole laps were traveled. In shuffling, distinguish a card label from its current position.
\Subhead{Satisfying stops / optional proof}
\textbf{K--1 stop:} compare one/two-hop routes and undo the hidden turn. \textbf{Optional:} explain missed shapes using separate loops.\par
\textbf{2--3 stop:} decode and compare +2/+3 routes. \textbf{Optional:} prove with alternate markings that +2 misses half, yet has an inverse.\par
\textbf{4--5 stop:} explain +4/+8 returns and loops on twelve places. \textbf{Optional:} prove the gcd formula with the adult strip scaffold; otherwise label the general rule a conjecture or supplied theorem.
\newpage
\PageTitle{FACILITATOR / 2}{Exact routes and the gcd rule}
\Subhead{K--1}
From circle, one hop visits triangle, square, star, circle (four turns); two hops alternate circle/square (two turns), never landing on triangle or star. A secret circle-to-square shift is two hops, so triangle becomes star, square becomes circle, and star becomes triangle. Two hops undo themselves. Three-hop turns follow circle, star, square, triangle, circle, visiting all four. On six positions, among steps 1 through 5, only 1 and 5 visit all six; 2 and 4 give two loops of length 3; 3 gives three loops of length 2.
\Subhead{Middle}
A to D is +3 on A--J. GDG decodes to \textbf{DAD}; BAG encodes to \textbf{EDJ}. B becomes E, so the additional B-to-F clue contradicts the fixed-shift rule. Repeated +2 gives A,C,E,G,I,A and B,D,F,H,J,B. Marking alternate letters explains why the two routes cannot meet: adding two preserves that marking. Repeated +3 gives A,D,G,J,C,F,I,B,E,H,A. The steps visiting all ten are \textbf{1,3,7,9}. Any shift is reversible: +2 is undone by -2, or +8. Skipping some letters along one repeated orbit is not information loss.
\Subhead{Upper: key and complete classification}
9 to 1 gives +4, which takes 2 to 6. Outputs 1,4,10 decode to \textbf{9,0,6}. Its four loops are (0,4,8), (1,5,9), (2,6,10), (3,7,11), each returning in three turns.
\begin{center}\renewcommand{\arraystretch}{1.25}\begin{tabular}{c|rrrrrrrrrrrr}Shift k&0&1&2&3&4&5&6&7&8&9&10&11\\\hline Return time&1&12&6&4&3&12&2&12&3&4&6&12\\Number of loops&12&1&2&3&4&1&6&1&4&3&2&1\end{tabular}\end{center}
Steps 1,5,7,11 visit all twelve. The zero rule is the identity and has first \emph{positive} return time 1, not zero. For eighteen positions, +6 returns after 3; +5 after 18.
\Subhead{Why the formula works}
For +8 on 12 places, return means \(3\mid2t\). Draw three t-strips, totaling 3t. If the first two strips total a multiple of 3, removing them leaves t a multiple of 3. Thus the first positive return is 3. This proves this case for every t. The optional general bridge is on page 4; do not treat a completed table as its proof.
\newpage
\PageTitle{FACILITATOR / 3}{From turning to shuffling}
\Subhead{Upper: inverse and composition}
On twelve places, +4 is undone by +8 (or -4). For every starting x, +4 followed by +7 gives x+11 modulo12; reversing the order gives the same because integer addition commutes. Every shift is bijective, even if its orbit from 0 is a proper subset. The visit-everywhere condition is \(\gcd(n,k)=1\), not a condition for reversibility. This distinction is central: a substitution permutation can have several cycles.
\Subhead{Extra: eight cards, three shuffles}
Rows, left to right with the leftmost card the top:
\begin{center}\renewcommand{\arraystretch}{1.5}\begin{tabular}{c|l}Shuffle&Row\\\hline 0&0 1 2 3 4 5 6 7\\1&0 4 1 5 2 6 3 7\\2&0 2 4 6 1 3 5 7\\3&0 1 2 3 4 5 6 7\end{tabular}\end{center}
The old-to-new position map is \textbf{0,2,4,6,1,3,5,7}. Its cycles are fixed 0, fixed 7, (1,2,4), and (3,6,5). Thus the whole shuffle returns after 3 and not earlier. This is also a direct week 3 loop argument.
For old position x in the first half (0--3), the new position is 2x. In the second half (4--7), it is \(2(x-4)+1=2x-7\). For x below 7 these are exactly \(2x\bmod7\). Position 7 stays 7, separately; it must not be merged with residue 0. Repetition gives \(2^t x\bmod7\). The powers 2,4,8 have remainders 2,4,1, so the first identity occurs at t=3. Tracking position 1 rules out earlier global returns, and remainder 1 fixes all positions 0--6.
\Subhead{Extra predictions}
Six cards use modulus 5. Powers 2,4,8,16 give remainders 2,4,3,1, so the order is \textbf{4}. Sixteen cards use modulus 15: 2,4,8,16 give 2,4,8,1, again order \textbf{4}. In general an out-shuffle on 2m cards maps positions below \(2m-1\) by doubling modulo \(2m-1\), fixes the final position, and has the multiplicative order of 2 as return time (the two-card identity is a trivial edge case). More cards can therefore have the same or even a smaller return time.
\Subhead{A useful additional contrast}
Random shuffling asks about a distribution over possible orders. This perfect shuffle is one precisely defined permutation repeated deterministically. Its return time says nothing by itself about how mixed a random shuffle is. Do not introduce chance unless it becomes a separate, explicitly changed rule.
\newpage
\PageTitle{FACILITATOR / 4}{Mathematical lineage and reuse}
\Subhead{Optional adult bridge: coprime divisibility}
If positive a,b are coprime and \(a\mid bt\), compare base strips a,b with scaled strips at,bt. Repeatedly subtract the smaller base strip from the larger, doing the same to the scaled strips. Both scaled lengths stay multiples of a. The positive base sum decreases, so the process ends with equal lengths c. Common divisors are unchanged: a divisor of x,y also divides x-y,y, and conversely by adding. Thus c=1, since the original lengths were coprime. The scaled process ends at t, which is therefore a multiple of a. Conversely, \(a\mid t\) plainly implies \(a\mid bt\).
\par
Try 5,3 first: 5-3=2, then 3-2=1. On scaled strips this leaves \(3t-(5t-3t)=t\). Equivalently, \(2\cdot3t-5t=t\); if \(5\mid3t\), the remainder is a multiple of 5. Ask the child to explain each removal. This is additional readiness, not assumed knowledge of prime factors.
\par
For k>0, set \(d=\gcd(n,k)\), \(n=da\), \(k=db\). Return means \(a\mid bt\); the lemma gives \(a\mid t\), hence first return n/d. Every start has this length, so there are d loops. For k=0, every point is fixed. If only the 3-and-2 case was explained, record that proof and label the general rule conjectured or given as a theorem. Full strip prompts and the Week 9 application: \texttt{plans/coprime-divisibility-scaffold.md}.
\Subhead{Sources and boundaries}
\textbf{Judson, \emph{Abstract Algebra: Theory and Applications}}, Ch. 4, \S4.1, Theorem 4.13, pp. 48--49, gives the cyclic order formula: \href{https://people.hsc.edu/faculty-staff/blins/books/JudsonAbstractAlgebra.pdf}{textbook}. Our rings are its additive version.
\textbf{Diaconis, Graham, Kantor (1983)}, \emph{The mathematics of perfect shuffles}, \S2, Lemma 1, p. 177, proves the out-shuffle order via doubling: \href{https://pages.uoregon.edu/kantor/PAPERS/PerfectShuffles.pdf}{original paper}. The extra proves this small return-time case and position rule, not the paper's shuffle-group classification. The plan has full research references.
Rozhkovskaya, Lesson 2, problems 2.3--2.7, supplies the shift-cipher entry. \emph{Math Circle by the Bay}, preface viii--x, supports manipulatives, changing pace, and reserves; Rozhkovskaya Lesson 8 records copying slowing the activity. Our rings, proof scaffolds, and hour are adaptations. Lowell Handout 6, problem 6.7, and Handout 7, problem 7.8, use a different card-ordering solitaire. Ask about recognition; do not assume this shuffle was taught.
\Subhead{Use history and checks}
Record actual steps, ring sizes, and extensions used as F04-K/M/U/X-v2. Keep 26-letter Caesar messages as a later alternative, not extra arithmetic for early finishers. Future years can compare prime/composite ring sizes, affine maps, in-shuffles, or perfect-shuffle groups; avoid replaying the same small cases.
\texttt{verify.py} checks all 12 shifts, the 10-letter routes, message examples, the 18-ring predictions, shuffle rows, position formulas, and first return times for 6, 8, 16 cards. The facilitator proofs establish why the patterns hold.

\end{document}
